% AIAA Journal of Aircraft -- full research article -- EARLY DRAFT
%
% Build:  pdflatex paper && bibtex paper && pdflatex paper && pdflatex paper
% with TEXINPUTS=./style: BSTINPUTS=./style:  (see README.md)
%
% STATUS: skeleton draft. Sections I-IV are written against implemented,
% tested code (Solver/DiskInduction.h, TestDiskInduction) and against the
% sweeps recorded in README.md. Section V is a stub: it needs the two-way
% coupling and the transition operating line, neither of which exists yet,
% and it must not be written from the one-way numbers as though they were
% final. Every quantitative claim below is traceable to a sweep listed in
% README.md -- do not add one that is not.

\documentclass[submit]{aiaa-tc}% journal mode: 12pt, double-spaced

\usepackage[utf8]{inputenc}
\usepackage{amsmath}
\usepackage{amssymb}
\usepackage{graphicx}
\usepackage{url}

\title{Aft-Fan Inflow Induction and Wing Separation\\
       During Tail-Sitter Transition}

% Same author block and open items as the companion drafts.
\author{
 Onur Tuncer%
  \thanks{Professor, Department of Aeronautical Engineering;
    onur.tuncer@itu.edu.tr. TODO: AIAA member grade.}\\
 {\normalsize\itshape
  Istanbul Technical University, Istanbul, T\"urkiye}
 \and
 \c{C}a\u{g}lar \"U\c{c}ler%
  \thanks{Associate Professor; caglar.ucler@ozyegin.edu.tr.
    TODO: department and AIAA member grade.}\\
 {\normalsize\itshape
  \"Ozye\u{g}in University, Istanbul, T\"urkiye}
 \and
 Ahmet G\"une\c{s}%
  \thanks{Associate Professor, Faculty of TODO; TODO: email and AIAA
    member grade.}\\
 {\normalsize\itshape
  Istanbul Technical University, Istanbul, T\"urkiye}
}

\newcommand{\eqnref}[1]{(\ref{#1})}

\begin{document}

\maketitle

\begin{abstract}
A tail-sitter with an aft-mounted ducted fan spends its transition at
incidences where the wing is separated, and the fan sits close enough
behind the wing to change that. The mechanism is not the one a blown-wing
study would model: the duct leading edge lies a fifth of a chord
\emph{aft} of the wing trailing edge, so its slipstream never reaches the
wing. What reaches the wing is the fan's \emph{upstream induction} --- the
accelerating field ahead of a loaded disk --- which is a favourable
streamwise pressure gradient and therefore delays separation. This paper
models that induction as the semi-infinite cylindrical vortex sheet
equivalent to a uniformly loaded actuator disk, verified against the
closed-form axial solution and against momentum theory, and applies it as
an external field to a coupled wing--body panel solve whose separation
behaviour is established independently. On the demonstrator airframe the
fan is found to buy $\Delta\alpha \approx 0.8^\circ$ of incidence at the
transition point, measured as the shift of the separation-location curve
rather than as the crossing of a single threshold --- a distinction that
matters, since a threshold evaluated on a two-degree grid reports no
effect at all from data containing this one. The benefit is concentrated
inboard, over precisely the span station that separates first, and
reverses sign outboard where the wing leaves the fan's streamtube and
meets its return flow.
\end{abstract}

\section*{Nomenclature}

\begin{tabbing}
  XXXXXXX \= \kill
  $A$ \> disk area, m$^2$ \\
  $b$ \> wing span, m \\
  $c_n$ \> chord in the leading-edge-normal plane, m \\
  $C_L$ \> lift coefficient \\
  $R$ \> disk radius, m \\
  $s$ \> axial distance from the disk plane, m \\
  $T$ \> fan thrust, N \\
  $u_x$ \> axial induced velocity, m/s \\
  $v_h$ \> hover induced velocity, $\sqrt{T/(2\rho A)}$, m/s \\
  $v_i$ \> induced velocity at the disk plane, m/s \\
  $V_\infty$ \> freestream speed, m/s \\
  $x_{\text{sep}}$ \> chordwise separation station, fraction of $c_n$ \\
  $\alpha$ \> angle of attack, deg \\
  $\gamma_t$ \> cylindrical vortex sheet strength, m/s \\
  $\mu$ \> velocity ratio, $V_\infty / v_h$ \\
  $\rho$ \> density, kg/m$^3$ \\
\end{tabbing}

\section{Introduction}

Transition is the part of a tail-sitter's envelope that sizes its wing and
its control authority, and it is the part a potential-flow method is least
able to speak to, because the wing is separated over much of it. The
companion study of this airframe establishes where separation begins and
how far forward it moves with incidence. This paper asks what the
propulsion system does about it.

The question is usually posed for a \emph{tractor} installation, where the
propeller slipstream washes the wing and raises its dynamic pressure --- a
blown wing, extensively studied. The airframe here is the other
arrangement, and the difference is not a detail. Its ducted fan sits at
the tail: the duct leading edge is $0.21$ chord \emph{behind} the wing
trailing edge, so the slipstream leaves downstream and never touches the
wing. A model that represents only the wake therefore reports exactly zero
interaction, which is a statement about the model rather than about the
aircraft.

What does reach the wing is the induction \emph{ahead} of the disk. A
loaded actuator disk accelerates the flow approaching it, from zero far
upstream to $v_i$ at the disk plane; a wing sitting in that field
experiences a favourable streamwise pressure gradient, and a favourable
gradient delays separation. The effect is the same in sign as a blown
wing's and entirely different in mechanism, and it is much less studied.

Two features of this configuration make the question sharp rather than
academic. First, the induced velocity is not a small correction: with the
annular disk area of the demonstrator, $v_h$ runs $14$--$24$~m/s against a
$25$~m/s cruise, so the induction is comparable to the flight speed
throughout transition and largest exactly where the airspeed is lowest.
Second, the geometry aligns: the duct outer radius is $0.21$ semi-span,
and the companion study places the wing's most advanced separation at
$|2y/b| \approx 0.15$--$0.23$. The part of the wing that sheds first is
the part the fan protects.

\section{The induction model}
\label{sec:model}

\subsection{Why the wake model cannot serve}

The solver's existing slipstream model reconstructs a rotor's wake from
annular momentum theory and returns zero induced velocity at any point
ahead of the disk plane. That is a correct economy for a surface lying in
the wake, and it is the wrong tool here by construction. Applied to a wing
upstream of the disk it returns zero, and the zero is the model's rather
than the flow's.

\subsection{The disk as a vortex cylinder}

A uniformly loaded actuator disk is exactly equivalent to a semi-infinite
cylindrical vortex sheet: azimuthal vorticity of constant strength
$\gamma_t = 2 v_i$ per unit length, shed from the disk edge and trailing
downstream. The equivalence is what makes the upstream field tractable,
because it turns ``what does a disk do ahead of itself'' into a
Biot--Savart integral over a known geometry with a closed-form answer on
the axis:
%
\begin{equation}
  \label{eq:axis}
  u_x(s) = \frac{\gamma_t}{2}\left[1 + \frac{s}{\sqrt{s^2 + R^2}}\right],
\end{equation}
%
with $s$ measured downstream from the disk plane. Equation~\eqnref{eq:axis}
returns $v_i$ at the disk, $2 v_i$ far downstream and zero far upstream ---
momentum theory's three values, recovered from the vortex system rather
than assumed --- and $v_i(1 - 1/\sqrt{2}) = 0.293\,v_i$ one radius
upstream, which is the number the whole interaction argument turns on.

Off the axis the sheet is discretized into vortex rings, each polygonized
into straight filaments and evaluated with the same Biot--Savart kernel
the lattice uses for its own wake. Ring spacing grows geometrically
downstream, since the first radius or two sets essentially the entire
upstream answer while the far tail varies slowly, and that spacing keys to
each cylinder's own radius rather than the disk's outer one.

A ducted fan around a tail boom is an \emph{annulus}, not a disk. A
uniformly loaded annulus carries no trailing vorticity between its radii,
so it sheds only at the two edges --- the outer at $+\gamma_t$, the inner
at $-\gamma_t$ --- and is therefore two superposed cylinders. On the axis
the two cancel at the disk plane and asymptotically but not in between:
downstream inside the bore the inner sheet dominates and the axial
induction reverses, which is the centerbody wake and a real feature of
annular jets.

Swirl is deliberately absent. The root vortex and the wake's axial
vorticity lie downstream of the disk and contribute nothing ahead of it,
which is the region this model exists to serve.

\subsection{Verification}

The discretized sheet is checked against Eq.~\eqnref{eq:axis} at stations
from four radii upstream to ten downstream, holding to 2\% of $v_i$, with
the transverse components vanishing on the axis by symmetry and the error
falling under refinement. The annulus is checked against its own closed
form, which pins the inner sheet's sign and strength, including the
reversed induction inside the bore. The momentum inversion
$T = 2\rho A v_i (V_\infty + v_i)$ round-trips, and forward speed reduces
$v_i$ at fixed thrust.

One further check is the one that magnitude comparisons cannot make: the
induction must point \emph{downstream} for a thrusting disk, both in the
jet and ahead of it, and must reverse when the disk axis reverses.
Inverting the ring circulation sense would turn the upstream acceleration
this paper rests on into a deceleration while leaving every magnitude
correct.

\section{Application to the demonstrator}
\label{sec:application}

The induction is applied as an external perturbation field to the coupled
wing--body--duct solve of the companion study, and the separation march is
run with the fan on and off at otherwise identical conditions. The
coupling is presently \emph{one-way}: the airframe sees the fan, the fan
does not see the airframe, so every result below is a lower bound on the
interaction.

\subsection{Magnitude and distribution}

The quantity that matters is not the induction at a point but its
variation \emph{along the chord}, because that variation is the pressure
gradient. The wing trailing edge sits roughly three times closer to the
disk than its leading edge, so the induction rises substantially across
the chord: at the transition point $V_\infty = 12$~m/s, $T = 25$~N, the
rise is $23\%$ of freestream at $2y/b = 0.10$, $14.5\%$ at $0.20$ and
$5\%$ at $0.30$. It is a favourable gradient, and it is strongest at the
trailing edge, which is where turbulent separation begins and from which
it marches forward.

Two features of the spanwise distribution deserve emphasis. It falls by a
factor of four between $2y/b = 0.10$ and $0.30$, concentrating the benefit
inboard --- over the stations that separate first. And beyond
$2y/b \approx 0.4$ it \emph{reverses}: outside the vortex cylinder the
induced axial velocity is a return flow, so the fan marginally penalizes
the outer wing. The penalty is under $2\%$ and the sign change is real,
not numerical, and it exists in the answer only because the upstream field
is modelled at all.

\subsection{How much incidence the fan buys}

% METHODOLOGICAL POINT -- keep this. The first pass of this study reported
% no effect, from data that contained a 0.8 degree shift.
Measuring the benefit requires care, and the obvious measurement fails.
Defining separation onset as a threshold crossing --- the first $\alpha$
at which any station separates forward of a fixed chordwise station ---
and evaluating it on the two-degree grid used for the coefficient tables
returns \emph{no change at all} between fan on and fan off. That result is
an artifact: a shift smaller than the grid spacing moves the crossing
within an interval without moving it across one, so the reported onset is
quantized to the grid and any sub-degree shift is invisible by
construction.

The appropriate measurement is the horizontal shift of the separation
curve $x_{\text{sep}}(\alpha)$, evaluated at several criterion levels.
Measured that way on a $0.5^\circ$ grid, the fan buys
$\Delta\alpha = +0.81^\circ$ averaged over criteria from
$x_{\text{sep}} = 0.85$ to $0.55$, with individual values between
$+0.73^\circ$ and $+0.86^\circ$. The consistency across criterion levels
is what makes this a property of the flow rather than of the definition;
the shift is smaller near onset, $+0.52^\circ$, where the separation point
lies at the trailing edge and the layer is least sensitive to the
gradient.

The benefit scales as expected with the velocity ratio: at cruise speed
the same comparison returns roughly a third of the transition value,
consistent with the induction scaling as $v_i/V_\infty$.

\section{Limitations}

The coupling is one-way, so the reported shift is a lower bound. The
induction model assumes uniform disk loading; real loading tapers at both
edges, softening the sheet into a band that superposed cylinders at
intermediate radii would represent. The fuselage boundary layer that an
aft fan genuinely ingests is not modelled, and it acts to reduce the
benefit. The analysis is quasi-steady, while transition is a manoeuvre
rather than a sequence of trim points. And the separation criterion
inherits the companion study's caveat: bubble bursting is not modelled, so
the separation boundary is an upper bound on usable incidence rather than
a predicted stall.

\section{Conclusions}
\label{sec:conclusions}

% STUB -- do not finalize. Needs the two-way coupling and the transition
% operating line (thrust and airspeed are currently independent parameters
% where a real transition walks a coupled schedule). See README.md.

\emph{[Draft: conclusions pending the two-way coupling and the transition
operating line. See the repository README.]}

\section*{Acknowledgments}

This work was supported by the Istanbul Technical University Scientific
Research Projects Coordination Unit (\.IT\"U BAP) under project
MGA-2026-48282, ``JETTAIL: Design, Control, and Autonomous Mission
Capability Development of a Jet-Flow-Vectored EDF-Based Tail-Sitter VTOL
UAV.''

% TODO: AI-assistance disclosure per AIAA policy, as in the companion
% drafts.

\bibliographystyle{aiaa}
\bibliography{references}

\end{document}
